\chapter{欧氏空间中的紧致性}
\label{chap:euclidean-compactness}

闭区间上的有限覆盖、收敛子列与连续函数的整体性质，已经在一维情形中
出现。欧氏空间保留这些结论，但证明需要把“区间二分”和“坐标有界”推广
到有限多个坐标。有限维性使坐标抽取在有限步后结束，也使一个有界集合能被
有限多个小立方体覆盖；这两件事正是 Heine--Borel 定理在 \(\R^n\) 中成立
的原因。

本章中的紧致仍以开覆盖的有限子覆盖性质为定义。除非特别说明，所有开、闭
和连续均相对于欧氏拓扑。

\section{多维 Bolzano--Weierstrass 定理}

\begin{theorem}[多维 Bolzano--Weierstrass 定理]
\label{thm:multidimensional-bolzano-weierstrass}
\(\R^n\) 中每个有界数列都有收敛子列。
\end{theorem}

\begin{proof}
写 \(x^{(k)}=(x_1^{(k)},\ldots,x_n^{(k)})\)。向量列有界意味着存在
\(M>0\)，使 \(\norm{x^{(k)}}_2\le M\)，故每个坐标列都有界。
先由一维 Bolzano--Weierstrass 定理，从原列抽出子列，使第一坐标收敛；
再从这个子列抽取使第二坐标收敛的子列。依次进行，经过 \(n\) 次后得到
子列 \(x^{(k_j)}\)，其全部坐标都收敛。由坐标收敛判据，这个子列在
\(\R^n\) 中收敛。
\end{proof}

这里不能为每个坐标各自选一个毫无关系的子列；第二次抽取必须在第一次所得
子列中进行，后续抽取才不会破坏已经得到的坐标收敛。因为坐标只有有限个，
这个过程能在有限步结束。

\begin{corollary}[有界无限集的聚点]
\label{cor:bounded-infinite-set-accumulation}
\(\R^n\) 中每个有界无限集都有聚点。
\end{corollary}

\begin{proof}
从集合中选取两两不同的点组成数列。由定理取收敛子列
\(x^{(k_j)}\to a\)。子列各项互异，所以 \(a\) 的每个删心邻域都含该
子列的某项，即 \(a\) 是原集合的聚点。
\end{proof}

若集合不有界，结论可能失败，例如 \(\Z\subset\R\) 没有聚点。反之，有界
条件只保证聚点存在于环境空间；要保证聚点仍属于原集合，还需要闭性。

\section{Heine--Borel 定理}

\begin{definition}[紧集]
集合 \(K\subset\R^n\) 称为紧集，如果它的每个开覆盖都有有限子覆盖。
\end{definition}

\subsection{有限网}

称有限点集
\(\{a_1,\ldots,a_N\}\subset\R^n\) 为 \(E\) 的一个
\(r\)-网，如果
\[
 E\subset\bigcup_{j=1}^N B(a_j,r).
\]

\begin{lemma}[欧氏有界集的有限网]
\label{lem:euclidean-bounded-finite-net}
若 \(E\subset\R^n\) 有界，则对每个 \(r>0\)，\(E\) 有有限
\(r\)-网。
\end{lemma}

\begin{proof}
取 \(M>0\) 使 \(E\subset[-M,M]^n\)。把每条坐标轴上的区间
\([-M,M]\) 等分成有限多个长度小于 \(r/\sqrt n\) 的小区间。所得有限
多个小立方体覆盖 \([-M,M]^n\)，每个与 \(E\) 相交的小立方体中选一点
\(a_j\in E\)。同一小立方体中任意两点的欧氏距离小于 \(r\)，故这些
\(a_j\) 构成有限 \(r\)-网。
\end{proof}

\subsection{闭有界刻画}

\begin{theorem}[欧氏空间的 Heine--Borel 定理]
\label{thm:euclidean-heine-borel}
集合 \(K\subset\R^n\) 紧，当且仅当它闭且有界。
\end{theorem}

\begin{proof}
若 \(K=\varnothing\)，闭性、有界性和紧致性均直接成立。以下设
\(K\ne\varnothing\)。

先设 \(K\) 紧。开球族 \(\{B(0,j):j\in\Npos\}\) 覆盖 \(K\)，有限
子覆盖包含于某个最大球 \(B(0,J)\)，所以 \(K\) 有界。若 \(a\notin K\)，
对每个 \(x\in K\) 取半径
\(r_x=\norm{x-a}_2/3>0\)。球 \(B(x,r_x)\) 覆盖 \(K\)，取有限子覆盖
对应于 \(x_1,\ldots,x_N\)，令 \(r=\min_jr_{x_j}>0\)。若
\(y\in B(a,r)\cap B(x_j,r_{x_j})\)，则
\[
 \norm{x_j-a}_2
 \le\norm{x_j-y}_2+\norm{y-a}_2
 <r_{x_j}+r\le2r_{x_j},
\]
与 \(\norm{x_j-a}_2=3r_{x_j}\) 矛盾。因此 \(B(a,r)\) 避开 \(K\)，
补集开，\(K\) 闭。

反过来，设 \(K\) 闭且有界，\(\mathcal U\) 是 \(K\) 的开覆盖。若它没有
有限子覆盖，则对每个 \(j\in\Npos\) 必存在 \(x_j\in K\)，使
\[
 B(x_j,1/j)\cap K
\]
不包含于任何单个 \(U\in\mathcal U\)。否则，若对某个 \(j\) 每一点都有
这样的覆盖成员，取 \(K\) 的有限 \(1/(2j)\)-网
\(a_1,\ldots,a_N\)。为每个网点选 \(U_i\in\mathcal U\)，使
\(B(a_i,1/j)\cap K\subset U_i\)；而每个 \(x\in K\) 距某个 \(a_i\)
小于 \(1/(2j)\)，故有限多个 \(U_i\) 已覆盖 \(K\)，矛盾。

数列 \((x_j)\) 有界，由多维 Bolzano--Weierstrass 定理存在子列
\(x_{j_k}\to x\)。闭性给出 \(x\in K\)。选 \(U\in\mathcal U\) 使
\(x\in U\)，再取 \(r>0\) 使 \(B(x,r)\subset U\)。充分大时
\(\norm{x_{j_k}-x}_2<r/2\) 且 \(1/j_k<r/2\)，从而
\[
 B(x_{j_k},1/j_k)\subset B(x,r)\subset U,
\]
与 \(x_{j_k}\) 的选择矛盾。因此原覆盖必有有限子覆盖。
\end{proof}

这条证明没有使用待证的紧致性来取得收敛子列：反向只调用上一节由一维
Bolzano--Weierstrass 定理建立的多维版本。证明同时给出欧氏紧集的一个
Lebesgue 数：对任意开覆盖，存在 \(\lambda>0\)，使每个 \(x\in K\) 的
相对球 \(B(x,\lambda)\cap K\) 包含于某个覆盖成员。有限网再把这个一致
尺度转成有限子覆盖。

\begin{theorem}[欧氏空间中的序列紧致性]
\label{thm:euclidean-sequential-compactness}
对 \(K\subset\R^n\)，以下条件等价：
\begin{enumerate}[label=\textup{(\roman*)}]
 \item \(K\) 紧；
 \item \(K\) 中每个数列都有收敛到 \(K\) 中某点的子列；
 \item \(K\) 闭且有界。
\end{enumerate}
\end{theorem}

\begin{proof}
上一条定理已给出 (i) 与 (iii) 的等价。若 (iii) 成立，任意 \(K\) 中
数列有界，由 Bolzano--Weierstrass 定理抽出收敛子列，闭性保证极限在
\(K\)。若 (ii) 成立而 \(K\) 无界，可取 \(\norm{x_j}_2>j\)，所得数列
无收敛子列；若 \(K\) 不闭，可取 \(x_j\in K\) 收敛到
\(x\notin K\)，其每个子列仍只能收敛到 \(x\)，也不可能满足 (ii)。
故 (ii) 推出 (iii)。
\end{proof}

\section{紧集上的连续函数}

\begin{theorem}[连续像保持紧致]
\label{thm:compact-continuous-image-euclidean}
若 \(K\subset\R^n\) 紧且 \(f:K\to\R^m\) 连续，则 \(f(K)\) 紧。
\end{theorem}

\begin{proof}
设 \(\{V_\lambda\}\) 是 \(f(K)\) 的开覆盖。连续性使
\(f^{-1}(V_\lambda)\) 成为 \(K\) 中的相对开集，并且这些原像覆盖
\(K\)。由紧致性选出有限个原像覆盖 \(K\)，相应的有限多个
\(V_\lambda\) 就覆盖 \(f(K)\)。
\end{proof}

\begin{corollary}[有界性与极值]
\label{cor:continuous-on-euclidean-compact}
紧集上的连续映射 \(f:K\to\R^m\) 有界。若 \(m=1\) 且
\(K\ne\varnothing\)，则 \(f\) 取得最大值和最小值。
\end{corollary}

\begin{proof}
像集 \(f(K)\) 紧，故由 Heine--Borel 定理有界。标量值时，\(f(K)\) 是
非空紧实数集，因而闭且有界；它的上确界和下确界属于 \(f(K)\)，分别由
某些 \(K\) 中点取到。
\end{proof}

\begin{theorem}[紧集上的一致连续性]
\label{thm:heine-cantor-euclidean}
若 \(K\subset\R^n\) 紧，\(f:K\to\R^m\) 连续，则 \(f\) 一致连续。
\end{theorem}

\begin{proof}
反设不一致连续。存在 \(\varepsilon_0>0\) 和点列 \(x_j,y_j\in K\)，使
\[
 \norm{x_j-y_j}_2<1/j,
 \qquad
 \norm{f(x_j)-f(y_j)}_2\ge\varepsilon_0.
\]
由序列紧致性，某子列 \(x_{j_k}\to x\in K\)。三角不等式说明
\(y_{j_k}\to x\)。连续性给出两条像子列都趋于 \(f(x)\)，故其距离趋零，
与下界矛盾。
\end{proof}

\begin{proposition}[到集合的距离]
\label{prop:distance-to-set-euclidean}
若 \(A\subset\R^n\) 非空，定义
\[
 d(x,A)=\inf_{a\in A}\norm{x-a}_2.
\]
则 \(x\mapsto d(x,A)\) 是 \(1\)-Lipschitz 函数。若 \(A\) 紧，则对每个
\(x\) 存在 \(a_x\in A\) 取得该距离。若非空集合 \(K,F\subset\R^n\)
分别为紧集和闭集，且 \(K\cap F=\varnothing\)，则
\[
 d(K,F):=\inf\{\norm{x-y}_2:x\in K,\ y\in F\}>0.
\]
\end{proposition}

\begin{proof}
对任意 \(a\in A\)，
\(d(x,A)\le\norm{x-a}_2\le\norm{x-y}_2+\norm{y-a}_2\)。对 \(a\)
取下确界得 \(d(x,A)\le\norm{x-y}_2+d(y,A)\)；交换 \(x,y\) 得
Lipschitz 估计。若 \(A\) 紧，连续函数
\(a\mapsto\norm{x-a}_2\) 取得最小值。

最后，按定义有 \(d(K,F)=\inf_{x\in K}d(x,F)\)。连续函数
\(x\mapsto d(x,F)\) 在紧集 \(K\) 上取得最小值。若最小值
为零，则相应 \(x_0\in K\) 满足 \(d(x_0,F)=0\)，即
\(x_0\in\overline F=F\)，与不相交矛盾。
\end{proof}

\section{连通、道路连通与凸性}

\begin{definition}[连通与道路连通]
集合 \(E\subset\R^n\) 称为连通的，如果不能写成两个非空、不交的相对
开集之并。若任意 \(x,y\in E\) 都存在连续映射
\(\gamma:[0,1]\to E\)，满足 \(\gamma(0)=x,\gamma(1)=y\)，则称
\(E\) 道路连通。
\end{definition}

道路连通推出连通。若 \(E=U\cup V\) 是分离，取 \(x\in U,y\in V\)
及连接它们的道路 \(\gamma\)。原像
\(\gamma^{-1}(U),\gamma^{-1}(V)\) 会把连通区间 \([0,1]\) 分离，矛盾。
反向在 \(\R^n\) 的任意子集上不成立；
\chapterref{chap:metric-connectedness}将给出标准反例。

\begin{definition}[凸集]
集合 \(C\subset\R^n\) 称为凸集，如果对任意 \(x,y\in C\) 与
\(0\le t\le1\)，都有
\[
 (1-t)x+ty\in C.
\]
\end{definition}

凸集道路连通，因为线段参数式
\(\gamma(t)=(1-t)x+ty\) 已给出道路。因此非空凸集必连通。开球、闭球、
线性子空间和半空间都是凸集；圆周连通且道路连通，却不是凸集。

\begin{theorem}[欧氏开集的连通性]
\label{thm:open-connected-path-connected-euclidean}
若 \(G\subset\R^n\) 开且连通，则 \(G\) 道路连通；更精确地，任意两点可由
有限条完全位于 \(G\) 中的线段首尾相接。
\end{theorem}

\begin{proof}
固定 \(x_0\in G\)，令 \(A\) 为能从 \(x_0\) 通过有限折线连接到的点集。
有 \(x_0\in A\)，故 \(A\ne\varnothing\)。若 \(x\in A\)，取
\(B(x,r)\subset G\)；球中任一点可先沿已有折线到 \(x\)，再沿一条线段
到达，故 \(B(x,r)\subset A\)，所以 \(A\) 在 \(G\) 中开。若
\(x\in G\setminus A\)，同样取 \(B(x,r)\subset G\)。若球中有一点属于
\(A\)，加上一条球内线段便可到达 \(x\)，矛盾；故该球包含于
\(G\setminus A\)，补集也相对开。连通性迫使 \(A=G\)。
\end{proof}

开性在这里不能删去；连通但非道路连通的集合通常正是在某个极限边缘处
失去局部开球。

\section{连续像和同胚}

\begin{theorem}[连续像保持连通]
\label{thm:continuous-image-connected-euclidean}
若 \(E\subset\R^n\) 连通且 \(f:E\to\R^m\) 连续，则 \(f(E)\) 连通。
若 \(E\) 道路连通，则 \(f(E)\) 道路连通。
\end{theorem}

\begin{proof}
若 \(f(E)=U\cup V\) 是分离，则两个非空原像
\(f^{-1}(U),f^{-1}(V)\) 会分离 \(E\)。道路情形更直接：连接 \(x,y\)
的道路与 \(f\) 复合，得到连接 \(f(x),f(y)\) 的道路。
\end{proof}

\begin{definition}[同胚]
连续双射 \(f:E\to F\) 若逆映射也连续，则称 \(f\) 为同胚，并称 \(E,F\)
同胚。紧致性、连通性与道路连通性在同胚下保持。
\end{definition}

连续双射的逆映射并不自动连续。例如
\(f:[0,2\pi)\to S^1\)、\(f(t)=(\cos t,\sin t)\) 是连续双射，但逆映射
在点 \((1,0)\) 处不连续。紧致定义域排除了这种“缺失端点被粘到一起”的
现象。

\begin{theorem}[紧空间到欧氏空间的连续双射]
\label{thm:compact-bijection-homeomorphism-euclidean}
若 \(K\subset\R^n\) 紧，\(f:K\to F\subset\R^m\) 是连续双射，则
\(f\) 是同胚。
\end{theorem}

\begin{proof}
若 \(C\subset K\) 相对闭，则 \(C\) 是紧集，连续像 \(f(C)\) 紧，从而在
\(\R^m\) 中闭，特别在 \(F\) 中相对闭。于是逆映射把闭集的原像写成
\[
 (f^{-1})^{-1}(C)=f(C),
\]
它是闭集。闭集原像判据给出 \(f^{-1}\) 连续。
\end{proof}

\section{闭有界与紧致：有限维性的作用}

Heine--Borel 定理中的“闭且有界”不能直接移到无限维函数空间。取
\(E=\Npos\)，在有界函数空间 \(\mathcal B(E)\) 中定义
\[
 e_k(j)=
 \begin{cases}
 1,&j=k,\\
 0,&j\ne k.
 \end{cases}
\]
则 \(\norm{e_k}_\infty=1\)，并且 \(k\ne l\) 时
\(\norm{e_k-e_l}_\infty=1\)。集合
\[
 A=\{e_k:k\in\Npos\}
\]
有界，也闭：任何在 \(A\) 中收敛的点列必须是 Cauchy 列，而由两两距离为
一可知它最终只能取同一个值。可是序列 \((e_k)\) 没有收敛子列，所以
\(A\) 不紧。

有限维证明依赖两项事实：有界列可以在有限次坐标抽取后同时收敛；有界集对
每个尺度都有有限网。上面的 \(A\) 在半径小于 \(1/2\) 时不可能有有限网，
因为一个这样的球至多含一个 \(e_k\)。
\chapterref{chap:compact-metric-spaces}会把“有限网”抽象成全有界性，
并给出度量空间中紧致性的准确替代判据。

\section*{习题}
\addcontentsline{toc}{section}{习题}

\subsection*{A组　子列与 Heine--Borel 定理}
\begin{enumerate}[label=\textbf{\thechapter.\arabic*.}]
 \exerciseitem{multidimensional-bw-diagonal}
 对 \(\R^3\) 中的有界列，逐次写出三次子列抽取及最终指标之间的包含关系。
 \exerciseitem{bounded-infinite-accumulation}
 从多维 Bolzano--Weierstrass 定理证明有界无限集有聚点，并说明选取两两
 不同点为何必要。
 \exerciseitem{euclidean-grid-net-count}
 若 \(E\subset[-M,M]^n\)，给出半径 \(r>0\) 的有限网点数的一个显式上界。
 \exerciseitem{compact-implies-closed-direct}
 不使用 Heine--Borel 定理的反向结论，按正文的球覆盖证明欧氏紧集闭。
 \exerciseitem{heine-borel-lebesgue-number}
 证明欧氏紧集的每个开覆盖都有 Lebesgue 数，并由有限网重新推出有限子覆盖。
 \exerciseitem{sequential-compactness-equivalence}
 完整证明欧氏空间中紧致、序列紧致与闭有界三者等价。
\end{enumerate}

\subsection*{B组　连续映射与几何结构}
\begin{enumerate}[label=\textbf{\thechapter.\arabic*.},resume]
 \exerciseitem{compact-product-euclidean}
 若 \(K\subset\R^n,L\subset\R^m\) 紧，证明
 \(K\times L\subset\R^{n+m}\) 紧。
 \exerciseitem{compact-sum-set}
 若 \(K,L\subset\R^n\) 紧，证明
 \(K+L=\{x+y:x\in K,y\in L\}\) 紧。
 \exerciseitem{compact-distance-attained}
 若非空紧集 \(K,L\subset\R^n\)，证明存在 \(x_0\in K,y_0\in L\) 使
 \[
  \norm{x_0-y_0}_2=\inf_{x\in K,y\in L}\norm{x-y}_2.
 \]
 \exerciseitem{compact-closed-positive-distance}
 证明若 \(K,F\subset\R^n\) 均非空，\(K\) 紧、\(F\) 闭且二者不相交，
 则 \(d(K,F)>0\)；再给出两个不相交
 闭集距离为零的欧氏空间例子。
 \exerciseitem{convex-connected}
 证明凸集道路连通；再给出道路连通但非凸的紧集。
 \exerciseitem{open-connected-polygonal}
 补全正文证明，并说明折线路径可以重新参数化为定义在 \([0,1]\) 上的连续
 道路。
 \exerciseitem{continuous-image-properties}
 分别证明连续像保持紧致、连通与道路连通，并比较三个证明使用的定义。
 \exerciseitem{compact-bijection-homeomorphism}
 证明紧集到欧氏子集的连续双射是同胚；用
 \([0,2\pi)\to S^1\) 说明紧致性不能删除。
\end{enumerate}

\subsection*{C组　有限维性与反例}
\begin{enumerate}[label=\textbf{\thechapter.\arabic*.},resume]
 \exerciseitem{unit-sphere-compact-any-norm}
 用有限维范数等价与 Heine--Borel 定理证明：\(\R^n\) 上任意范数的单位
 球面和闭单位球都紧。
 \exerciseitem{bounded-function-basis-closed}
 完整验证 \(\mathcal B(\Npos)\) 中集合
 \(\{e_k:k\in\Npos\}\) 闭、有界但不紧。
 \exerciseitem{continuous-proper-level-set}
 设 \(f:\R^n\to\R\) 连续且
 \(\norm{x}_2\to\infty\) 时 \(f(x)\to+\infty\)。证明 \(f\) 取得全局
 最小值。
 \optionalexerciseitem{cantor-set-compact-first-view}
 从 \([0,1]\) 开始每步删去各剩余区间的中间三分之一。证明所有阶段闭集的
 交是非空紧集，并证明其中每一点可由只含 \(0,2\) 的三进制展开表示。
 \optionalexerciseitem{space-filling-warning}
 说明“连续像不增加维数”作为朴素命题是错误的：查阅 Peano 曲线的准确
 陈述，并解释它为何不与连续像保持紧致和连通相冲突。本题只要求整理定理
 陈述和逻辑关系，不要求重建构造。
\end{enumerate}
